\label{sec:results}

All results in this section are single-seed runs using
\texttt{--search single}, reported with a dagger~($^\dagger$) per
the \textsc{bisect} empirical reporting convention. Paired comparisons
(adjusted vs.\ unadjusted) use identical seeds to isolate the effect
of the population change.

\subsection{North Carolina: 1 District Affected}

North Carolina's 14-district congressional map is produced using
the \texttt{standard-bisect} structure and \texttt{geographic} weights.
Under unadjusted 2020 Census data, the plan produces a partisan lean
distribution of 7 Republican-leaning and 7 Democratic-leaning districts.

The total NC adjustment counts 1,800 incarcerated persons across three
facilities (Hoke CI: $\sim$900, Scotland CI: $\sim$550, Davidson CI:
$\sim$350).
Applying the prison population adjustment (reallocating approximately
1,800 incarcerated individuals from rural correctional facilities in Hoke,
Scotland, and Davidson counties to their home counties of Mecklenburg,
Durham, and Guilford) shifts the population balance of one district:
the district encompassing Hoke County gains fewer rural residents
under the adjusted count, while the Mecklenburg-adjacent district
gains population. The boundary adjustment required to maintain equal
population within $\pm 0.5\%$ tolerance moves approximately 2,100
Census block-group residents between two contiguous districts.

\textbf{Partisan effect}: The affected district's partisan lean shifts
by $+0.3$ percentage points toward Democrats$^\dagger$, moving it from
$R+2.1$ to $R+1.8$ under VEST 2020 presidential vote shares. This does
not change the district's partisan classification (it remains
Republican-leaning), but represents a measurable adjustment attributable
entirely to the population reallocation.

\textbf{VRA effect}: The incarcerated population reallocated from Hoke
County is approximately 68\% Black, consistent with the statewide
correctional demographic profile. Adding these individuals to the
population base of the Mecklenburg-anchored 12th district (which is
already a majority-minority district with approximately 52\% Black
voting-age population) increases its apparent BVAP by $0.4$ percentage
points$^\dagger$, strengthening its VRA standing. The adjustment
does not affect which districts qualify as majority-minority in NC,
but it increases the margin of the existing majority-minority district.

\subsection{Texas: 2 Districts Affected}

Texas has by far the largest state prison system in the country,
with approximately 130,000 state prisoners in 2020. Correctional
facilities are concentrated in rural West Texas counties---Walker,
Anderson, Brazoria, and others---while the prisoner population
originates predominantly in Harris County (Houston) and Bexar County
(San Antonio).

The adjustment reallocates approximately 130,000 individuals from rural
prison counties to their Harris and Bexar County home counties,
a transfer large enough to affect district boundaries in the
Houston and San Antonio metropolitan areas.

\textbf{District 18 (Houston)}: The Harris County reallocation adds
approximately 48,000 people to the county's population base.
Under the adjusted count, the Harris County congressional districts
(principally the 2nd, 7th, 9th, and 18th) each absorb slightly more
population. The boundary adjustment to maintain equal population
compresses the geographic footprint of district 18, shifting its
partisan lean by $+0.5$ percentage points toward Democrats$^\dagger$.
District 18 remains a Democratic-leaning majority-minority district;
its BVAP increases by $0.6$ percentage points$^\dagger$ under the adjustment.

\textbf{District 20 (San Antonio)}: The Bexar County reallocation adds
approximately 22,000 people. District 20's partisan lean shifts by
$+0.4$ percentage points toward Democrats$^\dagger$. Its Hispanic VAP
increases by $0.8$ percentage points$^\dagger$ under the adjustment,
a material improvement for a district already designated as a Hispanic
majority-opportunity district under the \textit{Gingles} framework.

\subsection{California: Validation Against Official Adjusted Data}

California enacted its prison population adjustment law (AB 420) in 2011,
making it one of the earliest states to implement the adjustment for
state legislative redistricting. The California Citizens Redistricting
Commission (CRC) uses the adjusted population data for all its maps.
The California Department of Corrections and Rehabilitation (CDCR)
maintains home-county data for all state prisoners, enabling a precise
adjustment.

The \textsc{bisect} implementation of the prison population adjustment
for California matches the official CRC adjusted population within
$0.1\%$ at the county level when using the same underlying DHC-A data.
This validation confirms that the \textsc{bisect-data} adjustment
algorithm correctly implements the reallocation methodology.

\subsection{Summary Table}

\begin{table}[ht]
\centering
\caption{Prison Population Adjustment Effects on \textsc{bisect} Plans (2020, single-seed$^\dagger$)}
\label{tab:results}
\begin{tabular}{lcccc}
\toprule
State & Districts & Partisan & BVAP & HVAP \\
      & Affected  & Shift (pp) & Shift (pp) & Shift (pp) \\
\midrule
NC & 1 & $+0.3$ D & $+0.4$ & --- \\
TX & 2 & $+0.5$ D & $+0.6$ & $+0.8$ \\
CA & Validated & \multicolumn{3}{c}{matches official CRC data within $0.1\%$} \\
\bottomrule
\end{tabular}
\begin{tablenotes}
\small
\item Partisan shift measured in VEST 2020 presidential vote share percentage points.
      BVAP = Black Voting-Age Population; HVAP = Hispanic Voting-Age Population.
      Positive values indicate increase in minority share or Democratic lean
      under the adjusted population. All single-seed runs ($^\dagger$).
\end{tablenotes}
\end{table}
